\chapter{常见 AI 算子与模型结构}

理解 AI 算子之前，要先知道模型为什么由这些算子组成。推理引擎不是随意堆函数，而是在执行训练后确定下来的计算结构。不同模型结构使用不同算子，但许多底层模式相同：矩阵乘、逐元素变换、归一化、归约、查表、重排和数据搬运。

\section{线性层和 MLP}

线性层计算 $y = Wx + b$ 或批量形式 $Y = XW + b$。它把输入特征映射到新的特征空间，是 MLP、Transformer 投影层和分类头的基础。MLP 通常由线性层、激活函数和另一个线性层组成。LCQI 第一阶段选择 MLP，就是因为它用最小结构覆盖了参数、激活、矩阵运算和非线性。

激活函数让模型具有非线性表达能力。ReLU 是最简单的激活：小于零变成零，大于零保持不变。GELU、SiLU 等更常见于 Transformer，但计算更复杂。算子开发时，简单激活常常适合与前一个线性层输出融合，减少中间张量读写。

\section{卷积网络}

卷积网络常用于图像。卷积核在空间位置上滑动，提取局部模式。卷积算子涉及输入通道、输出通道、核大小、stride、padding 和布局。虽然卷积公式不同于矩阵乘，但实现上常常可以转化为 GEMM 或使用直接卷积、Winograd 等路径。读者理解卷积，是为了看到同一组硬件原则怎样作用在不同模型结构上。

\section{Embedding 和查表}

Embedding 把离散 token id 映射成向量。它本质上是查表：根据整数下标读取参数矩阵的一行。Embedding 的计算量不大，但内存访问可能不连续，容易受缓存和 TLB 影响。语言模型推理的第一步通常是 token id 查 Embedding，最后还可能用输出向量和词表矩阵计算 logits。

\section{Transformer}

Transformer 的核心部件包括多头 Attention、MLP、归一化、残差连接和位置编码。Attention 让每个 token 根据上下文聚合信息；MLP 对每个位置做特征变换；归一化稳定数值；残差连接帮助信息流动。推理时，prefill 阶段处理整段输入，decode 阶段逐 token 生成，并依赖 KV Cache 保存历史 key/value。

Transformer 看起来复杂，但落到底层仍然由算子组成：矩阵乘、Softmax、归一化、逐元素激活、查表、拼接和缓存读写。本册后续不会要求读者先成为模型训练专家，而是要能看懂这些算子在推理时怎样排列。

\section{算子接口}

一个算子接口至少要说明输入张量、输出张量、参数、shape 约束、dtype 约束和错误处理。算子内部可以有多种实现：标量版本用于正确性基线，SIMD 版本用于单核优化，多线程版本用于大规模输入，量化版本用于低带宽和低精度推理。接口稳定，内部实现才能逐步替换。

\begin{keyidea}
AI 模型结构决定算子列表，算子接口决定推理引擎边界，CPU 性能工程决定每个算子在本地 Linux 上能跑多快。
\end{keyidea}

\begin{exercisebox}
选择一个简单网络，例如两层 MLP 或小型 CNN，把它拆成算子列表。对每个算子写出输入 shape、输出 shape、参数 shape 和主要计算模式。不要写代码，先把推理图讲清楚。
\end{exercisebox}
